\documentclass{article}
\usepackage{PRIMEarxiv} % Core package for the PrimeAI Template
\usepackage{amsmath, amssymb, amsthm, bm, dcolumn} % Add amsthm here for the proof environment
\usepackage[numbers,sort&compress]{natbib} % Natbib for citations
\usepackage{graphicx} % For high-quality images
\usepackage{multicol} % Multi-column figures/tables
\usepackage{hyperref} % Hyperlinks
\usepackage{listings} % Code listings
\usepackage{authblk} % For structured affiliations
% Define theorem style (optional)
\theoremstyle{plain}
\newtheorem{theorem}{Theorem}
\renewcommand{\qedsymbol}{}

% Custom commands for primes and qubit encoding
\newcommand{\primeQubit}[1]{|p_{#1}\rangle}
\newcommand{\primeGate}[1]{U_{p_{#1}}}

\usepackage{wrapfig}
\usepackage[pscoord]{eso-pic}
\usepackage[fulladjust]{marginnote}
\reversemarginpar

% Typesetting improvements without footnote patching
\usepackage[protrusion=true, expansion=true, tracking=false]{microtype}
\microtypecontext{spacing=nonfrench}

% Line numbers
\usepackage[right]{lineno}

% Text layout - adjust as needed
\raggedright
\setlength{\parindent}{0.5cm}
\textwidth 5.25in 
\textheight 8.75in

% Set double spacing
\usepackage{setspace} 
\doublespacing

% Adjust width for specific content
\usepackage{changepage}

% Adjust caption style
\usepackage[aboveskip=1pt,labelfont=bf,labelsep=period,singlelinecheck=off]{caption}

% Remove brackets from references
\makeatletter
\renewcommand{\@biblabel}[1]{\quad#1.}
\makeatother

% Header, footer, and page numbers
\usepackage{lastpage,fancyhdr}
\pagestyle{fancy}
\fancyhf{}
\fancyhead[L]{Preprint - PrimeAI Enhanced Template}
\fancyfoot[C]{\scriptsize Multiplicity Theory © 2024 Ryan Van Gelder - Citizen Gardens \\ Licensed Under MIT and CC BY-NC-SA 4.0.}
\fancyfoot[R]{Page \thepage\ of \pageref{LastPage}}
\renewcommand{\footrule}{\hrule height 2pt \vspace{2mm}}

\begin{document}
\title{Revolutionizing Sourdough Bread \\ with Multiplicity Theory}

\author{Ryan Van Gelder}
\affil{Citizen Gardens - The Foundation of Multiplicity}
\date{\today}
\maketitle


\begin{abstract}
This paper presents a novel approach to optimizing sourdough fermentation using Multiplicity Theory, integrating prime-based encoding, quantum-inspired dynamics, and tensor interactions. By modeling ingredient interactions through quantum states and entanglement, this framework achieves precise control over fermentation variables, leading to superior flavor, texture, and microbial balance.
\end{abstract}

\section{Introduction}
Sourdough fermentation is a complex process driven by microbial interactions and environmental factors. Traditional methods rely on empirical knowledge, while this study applies Multiplicity Theory to encode and simulate the dynamics of fermentation, leveraging prime-based quantum models.

\section{Prime-Based Encoding of Ingredients}
Each ingredient is represented as a unique quantum state via prime encoding:\\
\[ \psi_{\text{sourdough}} = \sum_{i=1}^{n} c_i | p_i \rangle \]
where \( c_i \) represents the proportion of ingredient \( i \), and \( p_i \) is its assigned prime number:\\
\begin{align*}
\text{Flour: } p_1 = 2, \quad \text{Water: } p_2 = 3, \quad \text{Salt: } p_3 = 5, \quad \text{Starter: } p_4 = 7.
\end{align*}

\section{Microbial Entanglement}
The starter culture's microbial balance is modeled as an entangled quantum state:\\
\[
\psi_{\text{microbial}} = \frac{1}{\sqrt{2}} (| p_{\text{bacteria}} \rangle \otimes | p_{\text{yeast}} \rangle + | p_{\text{yeast}} \rangle \otimes | p_{\text{bacteria}} \rangle).
\]
This ensures harmony between lactic acid bacteria and yeast populations.

\section{Dynamic Feedback Mechanisms}
Environmental conditions such as temperature and hydration are regulated via the Dynamic Multiplicity Equation:\\
\[ \frac{\partial \rho_k}{\partial t} = \alpha_k \rho_k + \beta_k I_k + \gamma_k \sum_{j} T_{kj} \rho_j + \lambda (\Omega_B + \Omega_{FS}), \]
where:
\begin{itemize}
    \item \( \alpha_k \): Growth rate of microbial species.
    \item \( \beta_k \): Influence of hydration levels.
    \item \( \Omega_B, \Omega_{FS} \): Flavor and structure optimization parameters.
\end{itemize}

\section{Tensor Networks for Ingredient Interaction}
Interactions are captured through tensor networks:\\
\[ T_{ijk} = \sum_{l} \phi(p_{ijkl}) \]
where \( \phi(p_{ijkl}) \) models dependencies among hydration, gluten development, and microbial activity.

\section{Stochastic Enhancements}
Stochastic terms simulate variability, enriching texture and flavor:\\
\[ P(S, t) = \prod_{i=1}^{N} \left( p(x_i, t) + \epsilon_i(t) \right), \]
where \( \epsilon_i(t) \) introduces environmental noise.

\section{Simulation Framework}
\subsection{Quantum-Classical Synergy}
Classical systems preprocess ingredient mappings, while quantum frameworks simulate state evolution. Prime encoding modules adapt real-time feedback.

\subsection{Optimization Objectives}
The model aims to achieve:
\begin{enumerate}
    \item Enhanced microbial balance.
    \item Superior dough elasticity and flavor profiles.
    \item Reduced fermentation variability.
\end{enumerate}

\section{Conclusion}
Integrating Multiplicity Theory with sourdough fermentation establishes a groundbreaking paradigm for culinary science. Future work includes experimental validation and expanding applications to other fermented foods.



\nocite{*}
\bibliographystyle{plain}
\bibliography{references}

\end{document}